\documentclass[10pt]{article}
\usepackage[a4paper,margin=16mm]{geometry}
\usepackage{graphicx}
\usepackage{pgfplots}\pgfplotsset{compat=1.16}
\usepackage{quantikz}
\usepackage[field,figdir=.]{scqfig}
\usepackage[T1]{fontenc}\usepackage{helvet}\renewcommand{\familydefault}{\sfdefault}
\setlength{\parindent}{0pt}
\begin{document}
{\LARGE\bfseries scqfig gallery}\par\medskip
Every example in both styles: \textcolor{scqBlue}{\textbf{field}} (left) and \textbf{folio} (right). Build with \texttt{python examples/build\_gallery.py}; each panel comes from \texttt{examples/NN\_name.py}.\par\bigskip
\section*{\color{scqPurple}The LaTeX package}
\texttt{\textbackslash usepackage[field]\{scqfig\}} or \texttt{[folio]} defines the palette, both styles' roles (\texttt{scqFieldAccent}, \texttt{scqFolioAccent}) and the active style's roles under plain names (\texttt{scqAccent}). \texttt{\textbackslash scqgraphics\{t1\_decay\}} loads the figure of the active style, and the pgfplots styles \texttt{scqfig}, \texttt{scqfig data}, \texttt{scqfig model} and \texttt{scqfig ref} match the matplotlib output.\par\medskip
{\small
\begin{tabular}[t]{lcc}
\textbf{colour} & field & folio \\[2pt]
\texttt{scqInk} & \fcolorbox{scqLight}{scqFieldInk}{\rule{0pt}{7pt}\hspace{16pt}} & \fcolorbox{scqLight}{scqFolioInk}{\rule{0pt}{7pt}\hspace{16pt}} \\
\texttt{scqData} & \fcolorbox{scqLight}{scqFieldData}{\rule{0pt}{7pt}\hspace{16pt}} & \fcolorbox{scqLight}{scqFolioData}{\rule{0pt}{7pt}\hspace{16pt}} \\
\texttt{scqModel} & \fcolorbox{scqLight}{scqFieldModel}{\rule{0pt}{7pt}\hspace{16pt}} & \fcolorbox{scqLight}{scqFolioModel}{\rule{0pt}{7pt}\hspace{16pt}} \\
\texttt{scqSim} & \fcolorbox{scqLight}{scqFieldSim}{\rule{0pt}{7pt}\hspace{16pt}} & \fcolorbox{scqLight}{scqFolioSim}{\rule{0pt}{7pt}\hspace{16pt}} \\
\texttt{scqReference} & \fcolorbox{scqLight}{scqFieldReference}{\rule{0pt}{7pt}\hspace{16pt}} & \fcolorbox{scqLight}{scqFolioReference}{\rule{0pt}{7pt}\hspace{16pt}} \\
\texttt{scqAccent} & \fcolorbox{scqLight}{scqFieldAccent}{\rule{0pt}{7pt}\hspace{16pt}} & \fcolorbox{scqLight}{scqFolioAccent}{\rule{0pt}{7pt}\hspace{16pt}} \\
\texttt{scqSignal} & \fcolorbox{scqLight}{scqFieldSignal}{\rule{0pt}{7pt}\hspace{16pt}} & \fcolorbox{scqLight}{scqFolioSignal}{\rule{0pt}{7pt}\hspace{16pt}} \\
\texttt{scqBus} & \fcolorbox{scqLight}{scqFieldBus}{\rule{0pt}{7pt}\hspace{16pt}} & \fcolorbox{scqLight}{scqFolioBus}{\rule{0pt}{7pt}\hspace{16pt}} \\
\texttt{scqWash} & \fcolorbox{scqLight}{scqFieldWash}{\rule{0pt}{7pt}\hspace{16pt}} & \fcolorbox{scqLight}{scqFolioWash}{\rule{0pt}{7pt}\hspace{16pt}} \\
\texttt{scqOverlay} & \fcolorbox{scqLight}{scqFieldOverlay}{\rule{0pt}{7pt}\hspace{16pt}} & \fcolorbox{scqLight}{scqFolioOverlay}{\rule{0pt}{7pt}\hspace{16pt}} \\
\texttt{scqQubitA} & \fcolorbox{scqLight}{scqFieldQubitA}{\rule{0pt}{7pt}\hspace{16pt}} & \fcolorbox{scqLight}{scqFolioQubitA}{\rule{0pt}{7pt}\hspace{16pt}} \\
\texttt{scqQubitB} & \fcolorbox{scqLight}{scqFieldQubitB}{\rule{0pt}{7pt}\hspace{16pt}} & \fcolorbox{scqLight}{scqFolioQubitB}{\rule{0pt}{7pt}\hspace{16pt}} \\
\texttt{scqQubitC} & \fcolorbox{scqLight}{scqFieldQubitC}{\rule{0pt}{7pt}\hspace{16pt}} & \fcolorbox{scqLight}{scqFolioQubitC}{\rule{0pt}{7pt}\hspace{16pt}} \\
\end{tabular}
\hfill
\begin{tabular}[t]{lcc}
\textbf{colour} & field & folio \\[2pt]
\texttt{scqQubitD} & \fcolorbox{scqLight}{scqFieldQubitD}{\rule{0pt}{7pt}\hspace{16pt}} & \fcolorbox{scqLight}{scqFolioQubitD}{\rule{0pt}{7pt}\hspace{16pt}} \\
\texttt{scqLevel0} & \fcolorbox{scqLight}{scqFieldLevel0}{\rule{0pt}{7pt}\hspace{16pt}} & \fcolorbox{scqLight}{scqFolioLevel0}{\rule{0pt}{7pt}\hspace{16pt}} \\
\texttt{scqLevel1} & \fcolorbox{scqLight}{scqFieldLevel1}{\rule{0pt}{7pt}\hspace{16pt}} & \fcolorbox{scqLight}{scqFolioLevel1}{\rule{0pt}{7pt}\hspace{16pt}} \\
\texttt{scqLevel2} & \fcolorbox{scqLight}{scqFieldLevel2}{\rule{0pt}{7pt}\hspace{16pt}} & \fcolorbox{scqLight}{scqFolioLevel2}{\rule{0pt}{7pt}\hspace{16pt}} \\
\texttt{scqLevel3} & \fcolorbox{scqLight}{scqFieldLevel3}{\rule{0pt}{7pt}\hspace{16pt}} & \fcolorbox{scqLight}{scqFolioLevel3}{\rule{0pt}{7pt}\hspace{16pt}} \\
\texttt{scqTileNonlinear} & \fcolorbox{scqLight}{scqFieldTileNonlinear}{\rule{0pt}{7pt}\hspace{16pt}} & \fcolorbox{scqLight}{scqFolioTileNonlinear}{\rule{0pt}{7pt}\hspace{16pt}} \\
\texttt{scqTileCapacitive} & \fcolorbox{scqLight}{scqFieldTileCapacitive}{\rule{0pt}{7pt}\hspace{16pt}} & \fcolorbox{scqLight}{scqFolioTileCapacitive}{\rule{0pt}{7pt}\hspace{16pt}} \\
\texttt{scqTileInductive} & \fcolorbox{scqLight}{scqFieldTileInductive}{\rule{0pt}{7pt}\hspace{16pt}} & \fcolorbox{scqLight}{scqFolioTileInductive}{\rule{0pt}{7pt}\hspace{16pt}} \\
\texttt{scqChipPad} & \fcolorbox{scqLight}{scqFieldChipPad}{\rule{0pt}{7pt}\hspace{16pt}} & \fcolorbox{scqLight}{scqFolioChipPad}{\rule{0pt}{7pt}\hspace{16pt}} \\
\texttt{scqChipResonator} & \fcolorbox{scqLight}{scqFieldChipResonator}{\rule{0pt}{7pt}\hspace{16pt}} & \fcolorbox{scqLight}{scqFolioChipResonator}{\rule{0pt}{7pt}\hspace{16pt}} \\
\texttt{scqChipDrive} & \fcolorbox{scqLight}{scqFieldChipDrive}{\rule{0pt}{7pt}\hspace{16pt}} & \fcolorbox{scqLight}{scqFolioChipDrive}{\rule{0pt}{7pt}\hspace{16pt}} \\
\texttt{scqChipFlux} & \fcolorbox{scqLight}{scqFieldChipFlux}{\rule{0pt}{7pt}\hspace{16pt}} & \fcolorbox{scqLight}{scqFolioChipFlux}{\rule{0pt}{7pt}\hspace{16pt}} \\
\end{tabular}
}\par\bigskip
\noindent\begin{tikzpicture}\begin{axis}[scqfig field,width=7.4cm,height=4.6cm,xlabel={free evolution ($\mu$s)},ylabel={$P_1$},ymin=0,ymax=1,xmin=0,xmax=10,legend pos=north east]\addplot+[scqfig field data,samples=41,domain=0:10]{0.5+0.42*exp(-x/4)*cos(deg(2*pi*0.6*x))+0.02*sin(deg(37*x))};\addlegendentry{data}\addplot+[scqfig field model,samples=300,domain=0:10]{0.5+0.42*exp(-x/4)*cos(deg(2*pi*0.6*x))};\addlegendentry{fit}\end{axis}\end{tikzpicture}\hfill\begin{tikzpicture}\begin{axis}[scqfig folio,width=7.4cm,height=4.6cm,xlabel={free evolution ($\mu$s)},ylabel={$P_1$},ymin=0,ymax=1,xmin=0,xmax=10,legend pos=north east]\addplot+[scqfig folio data,samples=41,domain=0:10]{0.5+0.42*exp(-x/4)*cos(deg(2*pi*0.6*x))+0.02*sin(deg(37*x))};\addlegendentry{data}\addplot+[scqfig folio model,samples=300,domain=0:10]{0.5+0.42*exp(-x/4)*cos(deg(2*pi*0.6*x))};\addlegendentry{fit}\end{axis}\end{tikzpicture}\par
{\small pgfplots axes drawn natively by LaTeX with \texttt{[scqfig field]} (left) and \texttt{[scqfig folio]} (right).}\par\bigskip
\noindent\begin{quantikz}[scqfig circuit field,row sep=0.3cm,column sep=0.45cm]\lstick{$|0\rangle$} & \gate{H}\slice[label style=scq slice label]{$\Psi_1$} & \ctrl{1}\slice[label style=scq slice label]{$\Psi_2$} & \ctrl{2}\slice[label style=scq slice label]{$\Psi_3$} & \meter{}\slice[label style=scq slice label]{$\Psi_4$} & \\\lstick{$|0\rangle$} & \qw & \targ{} & \qw & \meter{} & \\\lstick{$|0\rangle$} & \qw & \qw & \targ{} & \meter{} &\end{quantikz}\hfill\begin{quantikz}[scqfig circuit folio,row sep=0.3cm,column sep=0.45cm]\lstick{$|0\rangle$} & \gate{H}\slice[label style=scq slice label]{$\Psi_1$} & \ctrl{1}\slice[label style=scq slice label]{$\Psi_2$} & \ctrl{2}\slice[label style=scq slice label]{$\Psi_3$} & \meter{}\slice[label style=scq slice label]{$\Psi_4$} & \\\lstick{$|0\rangle$} & \qw & \targ{} & \qw & \meter{} & \\\lstick{$|0\rangle$} & \qw & \qw & \targ{} & \meter{} &\end{quantikz}\par\smallskip
{\small The GHZ circuit of example 78 typeset by quantikz with \texttt{[scqfig circuit field]} (left) and \texttt{[scqfig circuit folio]} (right). Each slice tilts its label with \texttt{\textbackslash slice[label style=scq slice label]}.}\par\bigskip
\clearpage
\section*{\color{scqPurple}Spectrum}
\begin{minipage}{\textwidth}
\textbf{01\quad Potential well and wavefunctions}\hfill\texttt{01\_potential\_well.py}\par
{\small Transmon potential with the four lowest levels and their computed $|\psi(\varphi)|^2$ (charge basis, E\_J/E\_C = 80).}\par\smallskip
\begin{minipage}[t]{0.49\textwidth}\centering\includegraphics[width=\linewidth,height=0.36\textheight,keepaspectratio]{field/01_potential_well.pdf}\end{minipage}\hfill
\begin{minipage}[t]{0.49\textwidth}\centering\includegraphics[width=\linewidth,height=0.36\textheight,keepaspectratio]{folio/01_potential_well.pdf}\end{minipage}
\end{minipage}\par\bigskip\medskip
\begin{minipage}{\textwidth}
\textbf{02\quad Charge dispersion}\hfill\texttt{02\_charge\_dispersion.py}\par
{\small Charge dispersion of the three lowest levels as E\_J/E\_C grows from the Cooper-pair box to the transmon.}\par\smallskip
\begin{minipage}[t]{0.49\textwidth}\centering\includegraphics[width=\linewidth,height=0.36\textheight,keepaspectratio]{field/02_charge_dispersion.pdf}\end{minipage}\hfill
\begin{minipage}[t]{0.49\textwidth}\centering\includegraphics[width=\linewidth,height=0.36\textheight,keepaspectratio]{folio/02_charge_dispersion.pdf}\end{minipage}
\end{minipage}\par\bigskip\medskip
\begin{minipage}{\textwidth}
\textbf{03\quad Flux-tuned spectroscopy}\hfill\texttt{03\_flux\_spectroscopy.py}\par
{\small Two-tone spectroscopy of a SQUID transmon against flux, with the exact f01 and f02/2 drawn over the data (synthetic data).}\par\smallskip
\begin{minipage}[t]{0.49\textwidth}\centering\includegraphics[width=\linewidth,height=0.36\textheight,keepaspectratio]{field/03_flux_spectroscopy.pdf}\end{minipage}\hfill
\begin{minipage}[t]{0.49\textwidth}\centering\includegraphics[width=\linewidth,height=0.36\textheight,keepaspectratio]{folio/03_flux_spectroscopy.pdf}\end{minipage}
\end{minipage}\par\bigskip\medskip
\begin{minipage}{\textwidth}
\textbf{04\quad Avoided crossing}\hfill\texttt{04\_avoided\_crossing.py}\par
{\small Vacuum-Rabi splitting: resonator transmission as a tunable qubit crosses it, with the bare modes dashed and 2g marked (synthetic data).}\par\smallskip
\begin{minipage}[t]{0.49\textwidth}\centering\includegraphics[width=\linewidth,height=0.36\textheight,keepaspectratio]{field/04_avoided_crossing.pdf}\end{minipage}\hfill
\begin{minipage}[t]{0.49\textwidth}\centering\includegraphics[width=\linewidth,height=0.36\textheight,keepaspectratio]{folio/04_avoided_crossing.pdf}\end{minipage}
\end{minipage}\par\bigskip\medskip
\begin{minipage}{\textwidth}
\textbf{05\quad Scaling with E\_J/E\_C}\hfill\texttt{05\_ratio\_scaling.py}\par
{\small Transmon frequency, anharmonicity and charge dispersion against E\_J/E\_C, exact in the charge basis with the large-ratio asymptotes dashed and two designs marked (computed).}\par\smallskip
\begin{minipage}[t]{0.49\textwidth}\centering\includegraphics[width=\linewidth,height=0.36\textheight,keepaspectratio]{field/05_ratio_scaling.pdf}\end{minipage}\hfill
\begin{minipage}[t]{0.49\textwidth}\centering\includegraphics[width=\linewidth,height=0.36\textheight,keepaspectratio]{folio/05_ratio_scaling.pdf}\end{minipage}
\end{minipage}\par\bigskip\medskip
\begin{minipage}{\textwidth}
\textbf{06\quad Parity bands}\hfill\texttt{06\_parity\_bands.py}\par
{\small Even and odd charge-parity branches of a charge-sensitive transmon at E\_J/E\_C = 20: the 0-1 frequency with its splitting 2 chi\_0 and the zero-contrast offsets, and the first excited level (computed).}\par\smallskip
\begin{minipage}[t]{0.49\textwidth}\centering\includegraphics[width=\linewidth,height=0.36\textheight,keepaspectratio]{field/06_parity_bands.pdf}\end{minipage}\hfill
\begin{minipage}[t]{0.49\textwidth}\centering\includegraphics[width=\linewidth,height=0.36\textheight,keepaspectratio]{folio/06_parity_bands.pdf}\end{minipage}
\end{minipage}\par\bigskip\medskip
\begin{minipage}{\textwidth}
\textbf{07\quad Flux tuning}\hfill\texttt{07\_flux\_tuning.py}\par
{\small f01 of a SQUID transmon against flux for a symmetric SQUID and a 20 \% junction asymmetry, with the sweet spot and the flux that tunes it 1 GHz lower (computed).}\par\smallskip
\begin{minipage}[t]{0.49\textwidth}\centering\includegraphics[width=\linewidth,height=0.36\textheight,keepaspectratio]{field/07_flux_tuning.pdf}\end{minipage}\hfill
\begin{minipage}[t]{0.49\textwidth}\centering\includegraphics[width=\linewidth,height=0.36\textheight,keepaspectratio]{folio/07_flux_tuning.pdf}\end{minipage}
\end{minipage}\par\bigskip\medskip
\begin{minipage}{\textwidth}
\textbf{08\quad Dressed lines}\hfill\texttt{08\_dressed\_lines.py}\par
{\small Dressed one-excitation frequencies of a transmon (E\_C = 250 MHz) and a 7.5 GHz resonator coupled with g = 100 MHz, swept through the crossing, with the splitting and a design point marked (computed).}\par\smallskip
\begin{minipage}[t]{0.49\textwidth}\centering\includegraphics[width=\linewidth,height=0.36\textheight,keepaspectratio]{field/08_dressed_lines.pdf}\end{minipage}\hfill
\begin{minipage}[t]{0.49\textwidth}\centering\includegraphics[width=\linewidth,height=0.36\textheight,keepaspectratio]{folio/08_dressed_lines.pdf}\end{minipage}
\end{minipage}\par\bigskip\medskip
\section*{\color{scqPurple}Design}
\begin{minipage}{\textwidth}
\textbf{09\quad Junction design}\hfill\texttt{09\_junction\_curve.py}\par
{\small Qubit frequency against junction resistance for two transmons on an aluminium junction (gap 180 $\mu$eV), with L\_J on the top axis, a 4-6 GHz band, an allowed L\_J window and the designs with a 3 \% junction scatter (computed).}\par\smallskip
\begin{minipage}[t]{0.49\textwidth}\centering\includegraphics[width=\linewidth,height=0.36\textheight,keepaspectratio]{field/09_junction_curve.pdf}\end{minipage}\hfill
\begin{minipage}[t]{0.49\textwidth}\centering\includegraphics[width=\linewidth,height=0.36\textheight,keepaspectratio]{folio/09_junction_curve.pdf}\end{minipage}
\end{minipage}\par\bigskip\medskip
\begin{minipage}{\textwidth}
\textbf{10\quad Frequency plan}\hfill\texttt{10\_frequency\_plan.py}\par
{\small Frequency plan of a three-qubit chain: each transmon's f01 and f12, the readout resonators inside the Purcell-filter passband, and the collision windows each qubit's f01 must avoid (designed values).}\par\smallskip
\begin{minipage}[t]{0.49\textwidth}\centering\includegraphics[width=\linewidth,height=0.36\textheight,keepaspectratio]{field/10_frequency_plan.pdf}\end{minipage}\hfill
\begin{minipage}[t]{0.49\textwidth}\centering\includegraphics[width=\linewidth,height=0.36\textheight,keepaspectratio]{folio/10_frequency_plan.pdf}\end{minipage}
\end{minipage}\par\bigskip\medskip
\begin{minipage}{\textwidth}
\textbf{11\quad Collision map}\hfill\texttt{11\_collision\_map.py}\par
{\small Collision map of two coupled transmons: forbidden bands of the f\_q2 against f\_q1 plane, with three design points and their 20 MHz fabrication scatter; one point sits on a two-photon band (designed values).}\par\smallskip
\begin{minipage}[t]{0.49\textwidth}\centering\includegraphics[width=\linewidth,height=0.36\textheight,keepaspectratio]{field/11_collision_map.pdf}\end{minipage}\hfill
\begin{minipage}[t]{0.49\textwidth}\centering\includegraphics[width=\linewidth,height=0.36\textheight,keepaspectratio]{folio/11_collision_map.pdf}\end{minipage}
\end{minipage}\par\bigskip\medskip
\begin{minipage}{\textwidth}
\textbf{12\quad Capacitance matrix}\hfill\texttt{12\_capacitance\_matrix.py}\par
{\small Maxwell capacitance matrix of two transmons, their readout resonators and the feedline, coloured by \textbar{}C\textbar{} on a log scale (simulated values).}\par\smallskip
\begin{minipage}[t]{0.49\textwidth}\centering\includegraphics[width=\linewidth,height=0.36\textheight,keepaspectratio]{field/12_capacitance_matrix.pdf}\end{minipage}\hfill
\begin{minipage}[t]{0.49\textwidth}\centering\includegraphics[width=\linewidth,height=0.36\textheight,keepaspectratio]{folio/12_capacitance_matrix.pdf}\end{minipage}
\end{minipage}\par\bigskip\medskip
\begin{minipage}{\textwidth}
\textbf{13\quad Participation budget}\hfill\texttt{13\_participation.py}\par
{\small Surface-participation loss budget of a transmon: participation and loss tangent of each region, the quality factor each allows, and the total (simulated participation, assumed loss tangents).}\par\smallskip
\begin{minipage}[t]{0.49\textwidth}\centering\includegraphics[width=\linewidth,height=0.36\textheight,keepaspectratio]{field/13_participation.pdf}\end{minipage}\hfill
\begin{minipage}[t]{0.49\textwidth}\centering\includegraphics[width=\linewidth,height=0.36\textheight,keepaspectratio]{folio/13_participation.pdf}\end{minipage}
\end{minipage}\par\bigskip\medskip
\begin{minipage}{\textwidth}
\textbf{14\quad Dispersive shift}\hfill\texttt{14\_dispersive\_shift.py}\par
{\small Dispersive shift of a transmon against its frequency for a 7.15 GHz resonator, g = 85 MHz and alpha = -233 MHz, with the straddling regime shaded and the design point marked (computed).}\par\smallskip
\begin{minipage}[t]{0.49\textwidth}\centering\includegraphics[width=\linewidth,height=0.36\textheight,keepaspectratio]{field/14_dispersive_shift.pdf}\end{minipage}\hfill
\begin{minipage}[t]{0.49\textwidth}\centering\includegraphics[width=\linewidth,height=0.36\textheight,keepaspectratio]{folio/14_dispersive_shift.pdf}\end{minipage}
\end{minipage}\par\bigskip\medskip
\begin{minipage}{\textwidth}
\textbf{15\quad Purcell limit}\hfill\texttt{15\_purcell\_limit.py}\par
{\small Purcell limit on T1 against qubit frequency for a 7.15 GHz readout resonator (kappa = 1.5 MHz, g = 85 MHz), without and with a Q = 25 bandpass filter (computed).}\par\smallskip
\begin{minipage}[t]{0.49\textwidth}\centering\includegraphics[width=\linewidth,height=0.36\textheight,keepaspectratio]{field/15_purcell_limit.pdf}\end{minipage}\hfill
\begin{minipage}[t]{0.49\textwidth}\centering\includegraphics[width=\linewidth,height=0.36\textheight,keepaspectratio]{folio/15_purcell_limit.pdf}\end{minipage}
\end{minipage}\par\bigskip\medskip
\begin{minipage}{\textwidth}
\textbf{16\quad Convergence}\hfill\texttt{16\_convergence.py}\par
{\small Eigenmode frequency of a readout resonator against adaptive mesh pass, with the relative change per pass and a 0.02 \% target (synthetic data).}\par\smallskip
\begin{minipage}[t]{0.49\textwidth}\centering\includegraphics[width=\linewidth,height=0.36\textheight,keepaspectratio]{field/16_convergence.pdf}\end{minipage}\hfill
\begin{minipage}[t]{0.49\textwidth}\centering\includegraphics[width=\linewidth,height=0.36\textheight,keepaspectratio]{folio/16_convergence.pdf}\end{minipage}
\end{minipage}\par\bigskip\medskip
\begin{minipage}{\textwidth}
\textbf{17\quad Black-box quantization}\hfill\texttt{17\_black\_box.py}\par
{\small Black-box quantization of the transmon and readout resonator of example 72: (a) the susceptance the junction sees, whose rising zeros are the modes and whose slopes set their impedances; (b) the first-order Kerr matrix, with exact diagonalization of the same circuit in brackets (computed).}\par\smallskip
\begin{minipage}[t]{0.49\textwidth}\centering\includegraphics[width=\linewidth,height=0.36\textheight,keepaspectratio]{field/17_black_box.pdf}\end{minipage}\hfill
\begin{minipage}[t]{0.49\textwidth}\centering\includegraphics[width=\linewidth,height=0.36\textheight,keepaspectratio]{folio/17_black_box.pdf}\end{minipage}
\end{minipage}\par\bigskip\medskip
\begin{minipage}{\textwidth}
\textbf{18\quad Frequency lines}\hfill\texttt{18\_frequency\_lines.py}\par
{\small Frequency plan of a three-qubit sensor chip on one axis: f01 and f12 lines with their fabrication scatter, one flux scan, the collision windows of the tunable qubit, the readout resonators in a Purcell passband, the 20 mK thermal floor and the rest-energy axis (designed values).}\par\smallskip
\begin{minipage}[t]{0.49\textwidth}\centering\includegraphics[width=\linewidth,height=0.36\textheight,keepaspectratio]{field/18_frequency_lines.pdf}\end{minipage}\hfill
\begin{minipage}[t]{0.49\textwidth}\centering\includegraphics[width=\linewidth,height=0.36\textheight,keepaspectratio]{folio/18_frequency_lines.pdf}\end{minipage}
\end{minipage}\par\bigskip\medskip
\begin{minipage}{\textwidth}
\textbf{19\quad Design map}\hfill\texttt{19\_design\_map.py}\par
{\small Scan-rate gain R = eta\^{}2 dnu/Gamma of a photon counter over a quantum-limited amplifier on the plane of dark-count rate (over the signal linewidth) and efficiency, with R = 1 drawn heavier and three illustrative devices, one an efficiency lower bound (computed map, synthetic points).}\par\smallskip
\begin{minipage}[t]{0.49\textwidth}\centering\includegraphics[width=\linewidth,height=0.36\textheight,keepaspectratio]{field/19_design_map.pdf}\end{minipage}\hfill
\begin{minipage}[t]{0.49\textwidth}\centering\includegraphics[width=\linewidth,height=0.36\textheight,keepaspectratio]{folio/19_design_map.pdf}\end{minipage}
\end{minipage}\par\bigskip\medskip
\section*{\color{scqPurple}Readout}
\begin{minipage}{\textwidth}
\textbf{20\quad S21 of a resonator}\hfill\texttt{20\_s21\_resonator.py}\par
{\small Notch-resonator transmission as magnitude, phase and circle, with one complex fit and the cable delay, amplitude and phase calibrated out (synthetic data).}\par\smallskip
\begin{minipage}[t]{0.49\textwidth}\centering\includegraphics[width=\linewidth,height=0.36\textheight,keepaspectratio]{field/20_s21_resonator.pdf}\end{minipage}\hfill
\begin{minipage}[t]{0.49\textwidth}\centering\includegraphics[width=\linewidth,height=0.36\textheight,keepaspectratio]{folio/20_s21_resonator.pdf}\end{minipage}
\end{minipage}\par\bigskip\medskip
\begin{minipage}{\textwidth}
\textbf{21\quad Quality factor vs power}\hfill\texttt{21\_quality\_vs\_power.py}\par
{\small Internal quality factor of a resonator against mean photon number, with a fit of the saturable two-level-system model (synthetic data).}\par\smallskip
\begin{minipage}[t]{0.49\textwidth}\centering\includegraphics[width=\linewidth,height=0.36\textheight,keepaspectratio]{field/21_quality_vs_power.pdf}\end{minipage}\hfill
\begin{minipage}[t]{0.49\textwidth}\centering\includegraphics[width=\linewidth,height=0.36\textheight,keepaspectratio]{folio/21_quality_vs_power.pdf}\end{minipage}
\end{minipage}\par\bigskip\medskip
\begin{minipage}{\textwidth}
\textbf{22\quad Dispersive readout}\hfill\texttt{22\_dispersive\_readout.py}\par
{\small Dispersive readout: resonator transmission with the qubit in \textbar{}0\textgreater{} and in \textbar{}1\textgreater{}, the 2chi shift computed for g = 85 MHz and Delta = -2.09 GHz, and the readout frequency between the two dips (synthetic data).}\par\smallskip
\begin{minipage}[t]{0.49\textwidth}\centering\includegraphics[width=\linewidth,height=0.36\textheight,keepaspectratio]{field/22_dispersive_readout.pdf}\end{minipage}\hfill
\begin{minipage}[t]{0.49\textwidth}\centering\includegraphics[width=\linewidth,height=0.36\textheight,keepaspectratio]{folio/22_dispersive_readout.pdf}\end{minipage}
\end{minipage}\par\bigskip\medskip
\begin{minipage}{\textwidth}
\textbf{23\quad IQ blobs}\hfill\texttt{23\_iq\_blobs.py}\par
{\small Single-shot readout in the IQ plane: 4000 shots per prepared state with 1 and 2 sigma contours and the discrimination line; the \textbar{}1\textgreater{} shots include T1 decay during the 1.5 us integration (synthetic data).}\par\smallskip
\begin{minipage}[t]{0.49\textwidth}\centering\includegraphics[width=\linewidth,height=0.36\textheight,keepaspectratio]{field/23_iq_blobs.pdf}\end{minipage}\hfill
\begin{minipage}[t]{0.49\textwidth}\centering\includegraphics[width=\linewidth,height=0.36\textheight,keepaspectratio]{folio/23_iq_blobs.pdf}\end{minipage}
\end{minipage}\par\bigskip\medskip
\begin{minipage}{\textwidth}
\textbf{24\quad Readout histogram}\hfill\texttt{24\_readout\_histogram.py}\par
{\small Single-shot histograms projected on the discrimination axis, with Gaussian fits to each peak and the assignment fidelity; the tail between the peaks is T1 decay during readout (synthetic data).}\par\smallskip
\begin{minipage}[t]{0.49\textwidth}\centering\includegraphics[width=\linewidth,height=0.36\textheight,keepaspectratio]{field/24_readout_histogram.pdf}\end{minipage}\hfill
\begin{minipage}[t]{0.49\textwidth}\centering\includegraphics[width=\linewidth,height=0.36\textheight,keepaspectratio]{folio/24_readout_histogram.pdf}\end{minipage}
\end{minipage}\par\bigskip\medskip
\begin{minipage}{\textwidth}
\textbf{25\quad Assignment matrix}\hfill\texttt{25\_assignment\_matrix.py}\par
{\small Three-state assignment matrix from 5000 simulated shots per prepared state, assigned to the nearest IQ centre; decay during readout fills the lower triangle (synthetic data).}\par\smallskip
\begin{minipage}[t]{0.49\textwidth}\centering\includegraphics[width=\linewidth,height=0.36\textheight,keepaspectratio]{field/25_assignment_matrix.pdf}\end{minipage}\hfill
\begin{minipage}[t]{0.49\textwidth}\centering\includegraphics[width=\linewidth,height=0.36\textheight,keepaspectratio]{folio/25_assignment_matrix.pdf}\end{minipage}
\end{minipage}\par\bigskip\medskip
\section*{\color{scqPurple}Time domain}
\begin{minipage}{\textwidth}
\textbf{30\quad T1 decay}\hfill\texttt{30\_t1\_decay.py}\par
{\small Energy relaxation: excited-state population against delay with an exponential fit (synthetic data, T1 = 48 us).}\par\smallskip
\begin{minipage}[t]{0.49\textwidth}\centering\includegraphics[width=\linewidth,height=0.36\textheight,keepaspectratio]{field/30_t1_decay.pdf}\end{minipage}\hfill
\begin{minipage}[t]{0.49\textwidth}\centering\includegraphics[width=\linewidth,height=0.36\textheight,keepaspectratio]{folio/30_t1_decay.pdf}\end{minipage}
\end{minipage}\par\bigskip\medskip
\begin{minipage}{\textwidth}
\textbf{31\quad Power Rabi}\hfill\texttt{31\_rabi.py}\par
{\small Power Rabi: excited population against drive amplitude for a 40 ns pulse, with a cosine fit and the pi-pulse amplitude (synthetic data).}\par\smallskip
\begin{minipage}[t]{0.49\textwidth}\centering\includegraphics[width=\linewidth,height=0.36\textheight,keepaspectratio]{field/31_rabi.pdf}\end{minipage}\hfill
\begin{minipage}[t]{0.49\textwidth}\centering\includegraphics[width=\linewidth,height=0.36\textheight,keepaspectratio]{folio/31_rabi.pdf}\end{minipage}
\end{minipage}\par\bigskip\medskip
\begin{minipage}{\textwidth}
\textbf{32\quad Rabi chevron}\hfill\texttt{32\_chevron.py}\par
{\small Rabi chevron: excited population against drive detuning and pulse duration at a 10 MHz Rabi rate, with the first two population maxima of the model drawn over the left half (synthetic data).}\par\smallskip
\begin{minipage}[t]{0.49\textwidth}\centering\includegraphics[width=\linewidth,height=0.36\textheight,keepaspectratio]{field/32_chevron.pdf}\end{minipage}\hfill
\begin{minipage}[t]{0.49\textwidth}\centering\includegraphics[width=\linewidth,height=0.36\textheight,keepaspectratio]{folio/32_chevron.pdf}\end{minipage}
\end{minipage}\par\bigskip\medskip
\begin{minipage}{\textwidth}
\textbf{33\quad Ramsey fringes}\hfill\texttt{33\_ramsey.py}\par
{\small Ramsey fringes at 0.8 MHz detuning with a damped-cosine fit and its envelope (synthetic data, T2* = 4.2 us).}\par\smallskip
\begin{minipage}[t]{0.49\textwidth}\centering\includegraphics[width=\linewidth,height=0.36\textheight,keepaspectratio]{field/33_ramsey.pdf}\end{minipage}\hfill
\begin{minipage}[t]{0.49\textwidth}\centering\includegraphics[width=\linewidth,height=0.36\textheight,keepaspectratio]{folio/33_ramsey.pdf}\end{minipage}
\end{minipage}\par\bigskip\medskip
\begin{minipage}{\textwidth}
\textbf{34\quad Coherence decays}\hfill\texttt{34\_coherence\_decays.py}\par
{\small Ramsey envelope, Hahn echo and CPMG-16 decays of one qubit on one axis, each with its exponential fit and time constant (synthetic data).}\par\smallskip
\begin{minipage}[t]{0.49\textwidth}\centering\includegraphics[width=\linewidth,height=0.36\textheight,keepaspectratio]{field/34_coherence_decays.pdf}\end{minipage}\hfill
\begin{minipage}[t]{0.49\textwidth}\centering\includegraphics[width=\linewidth,height=0.36\textheight,keepaspectratio]{folio/34_coherence_decays.pdf}\end{minipage}
\end{minipage}\par\bigskip\medskip
\begin{minipage}{\textwidth}
\textbf{35\quad Randomized benchmarking}\hfill\texttt{35\_rb.py}\par
{\small Standard and interleaved randomized benchmarking of an X gate, 30 random sequences per length, with the error per Clifford and the interleaved gate error (synthetic data).}\par\smallskip
\begin{minipage}[t]{0.49\textwidth}\centering\includegraphics[width=\linewidth,height=0.36\textheight,keepaspectratio]{field/35_rb.pdf}\end{minipage}\hfill
\begin{minipage}[t]{0.49\textwidth}\centering\includegraphics[width=\linewidth,height=0.36\textheight,keepaspectratio]{folio/35_rb.pdf}\end{minipage}
\end{minipage}\par\bigskip\medskip
\begin{minipage}{\textwidth}
\textbf{36\quad Coherence over time}\hfill\texttt{36\_coherence\_timeseries.py}\par
{\small T1 and echo T2 measured every 10 minutes for 48 hours, with their histograms; T1 drops whenever a wandering defect drifts onto the qubit frequency (synthetic data).}\par\smallskip
\begin{minipage}[t]{0.49\textwidth}\centering\includegraphics[width=\linewidth,height=0.36\textheight,keepaspectratio]{field/36_coherence_timeseries.pdf}\end{minipage}\hfill
\begin{minipage}[t]{0.49\textwidth}\centering\includegraphics[width=\linewidth,height=0.36\textheight,keepaspectratio]{folio/36_coherence_timeseries.pdf}\end{minipage}
\end{minipage}\par\bigskip\medskip
\begin{minipage}{\textwidth}
\textbf{37\quad TLS spectral diffusion}\hfill\texttt{37\_tls\_map.py}\par
{\small T1 against qubit frequency over 24 hours; five defects show as Lorentzian T1 dips whose frequencies wander, and one switches between two frequencies (synthetic data).}\par\smallskip
\begin{minipage}[t]{0.49\textwidth}\centering\includegraphics[width=\linewidth,height=0.36\textheight,keepaspectratio]{field/37_tls_map.pdf}\end{minipage}\hfill
\begin{minipage}[t]{0.49\textwidth}\centering\includegraphics[width=\linewidth,height=0.36\textheight,keepaspectratio]{folio/37_tls_map.pdf}\end{minipage}
\end{minipage}\par\bigskip\medskip
\section*{\color{scqPurple}States}
\begin{minipage}{\textwidth}
\textbf{40\quad Bloch sphere}\hfill\texttt{40\_bloch\_sphere.py}\par
{\small Bloch sphere: (a) a detuned Rabi drive (Delta = 0.6 Omega) tilts the precession cone off the pole; (b) free precession from \textbar{}+\textgreater{} spirals to \textbar{}0\textgreater{} as T2 and T1 shrink the vector (computed).}\par\smallskip
\begin{minipage}[t]{0.49\textwidth}\centering\includegraphics[width=\linewidth,height=0.36\textheight,keepaspectratio]{field/40_bloch_sphere.pdf}\end{minipage}\hfill
\begin{minipage}[t]{0.49\textwidth}\centering\includegraphics[width=\linewidth,height=0.36\textheight,keepaspectratio]{folio/40_bloch_sphere.pdf}\end{minipage}
\end{minipage}\par\bigskip\medskip
\begin{minipage}{\textwidth}
\textbf{41\quad Wigner functions}\hfill\texttt{41\_wigner.py}\par
{\small Wigner functions of (a) the Fock state \textbar{}1\textgreater{} and (b) an even cat state with alpha = 2, on one colour scale of +-2/pi (computed).}\par\smallskip
\begin{minipage}[t]{0.49\textwidth}\centering\includegraphics[width=\linewidth,height=0.36\textheight,keepaspectratio]{field/41_wigner.pdf}\end{minipage}\hfill
\begin{minipage}[t]{0.49\textwidth}\centering\includegraphics[width=\linewidth,height=0.36\textheight,keepaspectratio]{folio/41_wigner.pdf}\end{minipage}
\end{minipage}\par\bigskip\medskip
\begin{minipage}{\textwidth}
\textbf{42\quad Density matrix}\hfill\texttt{42\_density\_matrix.py}\par
{\small Two-qubit tomography of the Bell state (\textbar{}00\textgreater{} + \textbar{}11\textgreater{})/sqrt(2): real and imaginary parts of the reconstructed density matrix (synthetic data).}\par\smallskip
\begin{minipage}[t]{0.49\textwidth}\centering\includegraphics[width=\linewidth,height=0.36\textheight,keepaspectratio]{field/42_density_matrix.pdf}\end{minipage}\hfill
\begin{minipage}[t]{0.49\textwidth}\centering\includegraphics[width=\linewidth,height=0.36\textheight,keepaspectratio]{folio/42_density_matrix.pdf}\end{minipage}
\end{minipage}\par\bigskip\medskip
\begin{minipage}{\textwidth}
\textbf{43\quad Density matrix in 3D}\hfill\texttt{43\_density\_bars3d.py}\par
{\small Real part of the reconstructed Bell-state density matrix as 3D bars, with the ideal state as wireframe outlines (synthetic data).}\par\smallskip
\begin{minipage}[t]{0.49\textwidth}\centering\includegraphics[width=\linewidth,height=0.36\textheight,keepaspectratio]{field/43_density_bars3d.pdf}\end{minipage}\hfill
\begin{minipage}[t]{0.49\textwidth}\centering\includegraphics[width=\linewidth,height=0.36\textheight,keepaspectratio]{folio/43_density_bars3d.pdf}\end{minipage}
\end{minipage}\par\bigskip\medskip
\section*{\color{scqPurple}Noise and errors}
\begin{minipage}{\textwidth}
\textbf{50\quad Noise spectral density}\hfill\texttt{50\_noise\_psd.py}\par
{\small Flux-noise spectral density with a power-law fit, amplitude at 1 Hz and a white floor (synthetic data).}\par\smallskip
\begin{minipage}[t]{0.49\textwidth}\centering\includegraphics[width=\linewidth,height=0.36\textheight,keepaspectratio]{field/50_noise_psd.pdf}\end{minipage}\hfill
\begin{minipage}[t]{0.49\textwidth}\centering\includegraphics[width=\linewidth,height=0.36\textheight,keepaspectratio]{folio/50_noise_psd.pdf}\end{minipage}
\end{minipage}\par\bigskip\medskip
\begin{minipage}{\textwidth}
\textbf{51\quad Error budget}\hfill\texttt{51\_error\_budget.py}\par
{\small CZ-gate error budget: contributions, their sum, and the measured error with the unexplained part hatched.}\par\smallskip
\begin{minipage}[t]{0.49\textwidth}\centering\includegraphics[width=\linewidth,height=0.36\textheight,keepaspectratio]{field/51_error_budget.pdf}\end{minipage}\hfill
\begin{minipage}[t]{0.49\textwidth}\centering\includegraphics[width=\linewidth,height=0.36\textheight,keepaspectratio]{folio/51_error_budget.pdf}\end{minipage}
\end{minipage}\par\bigskip\medskip
\begin{minipage}{\textwidth}
\textbf{52\quad Correlated error bursts}\hfill\texttt{52\_burst\_raster.py}\par
{\small Correlated error bursts: detections per qubit over time, with the flagged fraction and burst marks above (synthetic data from a toy burst model).}\par\smallskip
\begin{minipage}[t]{0.49\textwidth}\centering\includegraphics[width=\linewidth,height=0.36\textheight,keepaspectratio]{field/52_burst_raster.pdf}\end{minipage}\hfill
\begin{minipage}[t]{0.49\textwidth}\centering\includegraphics[width=\linewidth,height=0.36\textheight,keepaspectratio]{folio/52_burst_raster.pdf}\end{minipage}
\end{minipage}\par\bigskip\medskip
\begin{minipage}{\textwidth}
\textbf{53\quad Error correlations}\hfill\texttt{53\_error\_correlation.py}\par
{\small Pairwise correlation of error detections across a 24-qubit chain; bursts correlate neighbours (synthetic data from a toy burst model).}\par\smallskip
\begin{minipage}[t]{0.49\textwidth}\centering\includegraphics[width=\linewidth,height=0.36\textheight,keepaspectratio]{field/53_error_correlation.pdf}\end{minipage}\hfill
\begin{minipage}[t]{0.49\textwidth}\centering\includegraphics[width=\linewidth,height=0.36\textheight,keepaspectratio]{folio/53_error_correlation.pdf}\end{minipage}
\end{minipage}\par\bigskip\medskip
\section*{\color{scqPurple}Dark matter}
\begin{minipage}{\textwidth}
\textbf{55\quad Rotation curve}\hfill\texttt{55\_rotation\_curve.py}\par
{\small Rotation curve of a model spiral galaxy: the speed that the visible matter alone would give falls off outside the disk, while a dark halo keeps the total flat near 220 km/s (illustrative model, synthetic points).}\par\smallskip
\begin{minipage}[t]{0.49\textwidth}\centering\includegraphics[width=\linewidth,height=0.36\textheight,keepaspectratio]{field/55_rotation_curve.pdf}\end{minipage}\hfill
\begin{minipage}[t]{0.49\textwidth}\centering\includegraphics[width=\linewidth,height=0.36\textheight,keepaspectratio]{folio/55_rotation_curve.pdf}\end{minipage}
\end{minipage}\par\bigskip\medskip
\begin{minipage}{\textwidth}
\textbf{56\quad Occupation number}\hfill\texttt{56\_occupation\_number.py}\par
{\small Dark-matter particles per de Broglie volume against rest energy for 0.45 GeV/cm³ and 220 km/s: below about 40 eV each cell holds many particles and the dark matter acts as one classical wave, which covers the 1 to 10 GHz band of qubit searches (computed).}\par\smallskip
\begin{minipage}[t]{0.49\textwidth}\centering\includegraphics[width=\linewidth,height=0.36\textheight,keepaspectratio]{field/56_occupation_number.pdf}\end{minipage}\hfill
\begin{minipage}[t]{0.49\textwidth}\centering\includegraphics[width=\linewidth,height=0.36\textheight,keepaspectratio]{folio/56_occupation_number.pdf}\end{minipage}
\end{minipage}\par\bigskip\medskip
\begin{minipage}{\textwidth}
\textbf{57\quad Dark-matter line shape}\hfill\texttt{57\_dm\_lineshape.py}\par
{\small Line shape of a 5 GHz dark-matter wave for isotropic Maxwell speeds with v0 = 220 km/s: the line is a few kHz wide, one part in a million of its frequency, with a mean offset of 2.0 kHz (computed).}\par\smallskip
\begin{minipage}[t]{0.49\textwidth}\centering\includegraphics[width=\linewidth,height=0.36\textheight,keepaspectratio]{field/57_dm_lineshape.pdf}\end{minipage}\hfill
\begin{minipage}[t]{0.49\textwidth}\centering\includegraphics[width=\linewidth,height=0.36\textheight,keepaspectratio]{folio/57_dm_lineshape.pdf}\end{minipage}
\end{minipage}\par\bigskip\medskip
\begin{minipage}{\textwidth}
\textbf{58\quad Dark-matter field trace}\hfill\texttt{58\_dm\_field\_trace.py}\par
{\small Amplitude and phase of a 5 GHz dark-matter wave built from 400 tones with standard-halo velocities: both wander on the 200 $\mu$s coherence time, twice the length of the shaded 100 $\mu$s qubit shot (simulated, seeded).}\par\smallskip
\begin{minipage}[t]{0.49\textwidth}\centering\includegraphics[width=\linewidth,height=0.36\textheight,keepaspectratio]{field/58_dm_field_trace.pdf}\end{minipage}\hfill
\begin{minipage}[t]{0.49\textwidth}\centering\includegraphics[width=\linewidth,height=0.36\textheight,keepaspectratio]{folio/58_dm_field_trace.pdf}\end{minipage}
\end{minipage}\par\bigskip\medskip
\begin{minipage}{\textwidth}
\textbf{59\quad Dark-photon reach}\hfill\texttt{59\_reach.py}\par
{\small Kinetic-mixing reach of a qubit dark-photon search against frequency for one and sixteen sensors at 20 and 50 mK, from an illustrative counting model with thermal dark counts, an excluded region and a synthetic measured limit, with the dark-photon mass on the top axis (computed).}\par\smallskip
\begin{minipage}[t]{0.49\textwidth}\centering\includegraphics[width=\linewidth,height=0.36\textheight,keepaspectratio]{field/59_reach.pdf}\end{minipage}\hfill
\begin{minipage}[t]{0.49\textwidth}\centering\includegraphics[width=\linewidth,height=0.36\textheight,keepaspectratio]{folio/59_reach.pdf}\end{minipage}
\end{minipage}\par\bigskip\medskip
\section*{\color{scqPurple}Processor}
\begin{minipage}{\textwidth}
\textbf{60\quad Device map}\hfill\texttt{60\_device\_map.py}\par
{\small Device map of a 3 x 4 qubit grid with T1 on the qubits and CZ error on the couplers (synthetic data).}\par\smallskip
\begin{minipage}[t]{0.49\textwidth}\centering\includegraphics[width=\linewidth,height=0.36\textheight,keepaspectratio]{field/60_device_map.pdf}\end{minipage}\hfill
\begin{minipage}[t]{0.49\textwidth}\centering\includegraphics[width=\linewidth,height=0.36\textheight,keepaspectratio]{folio/60_device_map.pdf}\end{minipage}
\end{minipage}\par\bigskip\medskip
\begin{minipage}{\textwidth}
\textbf{61\quad Error distributions}\hfill\texttt{61\_error\_cdf.py}\par
{\small Cumulative distributions of single-qubit, CZ and readout errors across a 54-qubit processor, medians marked (synthetic data).}\par\smallskip
\begin{minipage}[t]{0.49\textwidth}\centering\includegraphics[width=\linewidth,height=0.36\textheight,keepaspectratio]{field/61_error_cdf.pdf}\end{minipage}\hfill
\begin{minipage}[t]{0.49\textwidth}\centering\includegraphics[width=\linewidth,height=0.36\textheight,keepaspectratio]{folio/61_error_cdf.pdf}\end{minipage}
\end{minipage}\par\bigskip\medskip
\section*{\color{scqPurple}Schematics}
\begin{minipage}{\textwidth}
\textbf{70\quad Chip plan}\hfill\texttt{70\_chip\_plan.py}\par
{\small Plan view of a two-qubit chip: Xmons joined by a tunable coupler, readout resonators on a feedline, drive and flux lines, airbridges.}\par\smallskip
\begin{minipage}[t]{0.49\textwidth}\centering\includegraphics[width=\linewidth,height=0.36\textheight,keepaspectratio]{field/70_chip_plan.pdf}\end{minipage}\hfill
\begin{minipage}[t]{0.49\textwidth}\centering\includegraphics[width=\linewidth,height=0.36\textheight,keepaspectratio]{folio/70_chip_plan.pdf}\end{minipage}
\end{minipage}\par\bigskip\medskip
\begin{minipage}{\textwidth}
\textbf{71\quad Chip on a slab}\hfill\texttt{71\_chip\_oblique.py}\par
{\small The two-qubit chip of example 70 laid on a slab by one rigid oblique map, with a plan-view zoom on the SQUID of Q1.}\par\smallskip
\begin{minipage}[t]{0.49\textwidth}\centering\includegraphics[width=\linewidth,height=0.36\textheight,keepaspectratio]{field/71_chip_oblique.pdf}\end{minipage}\hfill
\begin{minipage}[t]{0.49\textwidth}\centering\includegraphics[width=\linewidth,height=0.36\textheight,keepaspectratio]{folio/71_chip_oblique.pdf}\end{minipage}
\end{minipage}\par\bigskip\medskip
\begin{minipage}{\textwidth}
\textbf{72\quad Transmon and readout circuit}\hfill\texttt{72\_lumped\_circuit.py}\par
{\small Lumped circuit of a transmon coupled to a quarter-wave readout resonator on a feedline, each Hamiltonian term on the tile of the element that produces it.}\par\smallskip
\begin{minipage}[t]{0.49\textwidth}\centering\includegraphics[width=\linewidth,height=0.36\textheight,keepaspectratio]{field/72_lumped_circuit.pdf}\end{minipage}\hfill
\begin{minipage}[t]{0.49\textwidth}\centering\includegraphics[width=\linewidth,height=0.36\textheight,keepaspectratio]{folio/72_lumped_circuit.pdf}\end{minipage}
\end{minipage}\par\bigskip\medskip
\begin{minipage}{\textwidth}
\textbf{73\quad Cryostat wiring}\hfill\texttt{73\_cryostat\_wiring.py}\par
{\small Wiring of a dilution refrigerator: attenuated and filtered XY, Z and readout-input lines, and a readout output through isolators, a TWPA and a HEMT.}\par\smallskip
\begin{minipage}[t]{0.49\textwidth}\centering\includegraphics[width=\linewidth,height=0.36\textheight,keepaspectratio]{field/73_cryostat_wiring.pdf}\end{minipage}\hfill
\begin{minipage}[t]{0.49\textwidth}\centering\includegraphics[width=\linewidth,height=0.36\textheight,keepaspectratio]{folio/73_cryostat_wiring.pdf}\end{minipage}
\end{minipage}\par\bigskip\medskip
\begin{minipage}{\textwidth}
\textbf{74\quad Pulse sequence}\hfill\texttt{74\_pulse\_sequence.py}\par
{\small Pulse sequence of a conditional-phase measurement: DRAG pulses on two qubits, a CZ flux pulse and a dispersive readout with its integration window.}\par\smallskip
\begin{minipage}[t]{0.49\textwidth}\centering\includegraphics[width=\linewidth,height=0.36\textheight,keepaspectratio]{field/74_pulse_sequence.pdf}\end{minipage}\hfill
\begin{minipage}[t]{0.49\textwidth}\centering\includegraphics[width=\linewidth,height=0.36\textheight,keepaspectratio]{folio/74_pulse_sequence.pdf}\end{minipage}
\end{minipage}\par\bigskip\medskip
\begin{minipage}{\textwidth}
\textbf{75\quad Level ladder and transitions}\hfill\texttt{75\_level\_ladder.py}\par
{\small Transmon ladder (E\_J/E\_C = 40) to scale, with drive photons for f01, f12 and the two-photon f02/2, and the parity-forbidden single-photon 0-2 transition.}\par\smallskip
\begin{minipage}[t]{0.49\textwidth}\centering\includegraphics[width=\linewidth,height=0.36\textheight,keepaspectratio]{field/75_level_ladder.pdf}\end{minipage}\hfill
\begin{minipage}[t]{0.49\textwidth}\centering\includegraphics[width=\linewidth,height=0.36\textheight,keepaspectratio]{folio/75_level_ladder.pdf}\end{minipage}
\end{minipage}\par\bigskip\medskip
\begin{minipage}{\textwidth}
\textbf{76\quad Film stack}\hfill\texttt{76\_film\_stack.py}\par
{\small Cross-section of a tantalum CPW on silicon with its metal-air, metal-substrate and substrate-air interfaces marked.}\par\smallskip
\begin{minipage}[t]{0.49\textwidth}\centering\includegraphics[width=\linewidth,height=0.36\textheight,keepaspectratio]{field/76_film_stack.pdf}\end{minipage}\hfill
\begin{minipage}[t]{0.49\textwidth}\centering\includegraphics[width=\linewidth,height=0.36\textheight,keepaspectratio]{folio/76_film_stack.pdf}\end{minipage}
\end{minipage}\par\bigskip\medskip
\begin{minipage}{\textwidth}
\textbf{77\quad Qubit array in a uniform field}\hfill\texttt{77\_spin\_array.py}\par
{\small A 3 x 4 qubit array drawn as spins in a uniform field (the wash), with photons arriving from several directions and three qubits excited.}\par\smallskip
\begin{minipage}[t]{0.49\textwidth}\centering\includegraphics[width=\linewidth,height=0.36\textheight,keepaspectratio]{field/77_spin_array.pdf}\end{minipage}\hfill
\begin{minipage}[t]{0.49\textwidth}\centering\includegraphics[width=\linewidth,height=0.36\textheight,keepaspectratio]{folio/77_spin_array.pdf}\end{minipage}
\end{minipage}\par\bigskip\medskip
\begin{minipage}{\textwidth}
\textbf{78\quad Quantum circuit}\hfill\texttt{78\_quantum\_circuit.py}\par
{\small GHZ-state preparation: a Hadamard and two CNOTs entangle three qubits, then all three are measured; dashed slices mark the states Psi\_1 to Psi\_4.}\par\smallskip
\begin{minipage}[t]{0.49\textwidth}\centering\includegraphics[width=\linewidth,height=0.36\textheight,keepaspectratio]{field/78_quantum_circuit.pdf}\end{minipage}\hfill
\begin{minipage}[t]{0.49\textwidth}\centering\includegraphics[width=\linewidth,height=0.36\textheight,keepaspectratio]{folio/78_quantum_circuit.pdf}\end{minipage}
\end{minipage}\par\bigskip\medskip
\begin{minipage}{\textwidth}
\textbf{79\quad Teleportation and gate glyphs}\hfill\texttt{79\_teleportation.py}\par
{\small (a) Teleportation: a Bell pair, a Bell measurement and classically controlled X and Z corrections. (b) The other gate glyphs: rotation, controlled-U with an open control, CZ, SWAP and a three-qubit box.}\par\smallskip
\begin{minipage}[t]{0.49\textwidth}\centering\includegraphics[width=\linewidth,height=0.36\textheight,keepaspectratio]{field/79_teleportation.pdf}\end{minipage}\hfill
\begin{minipage}[t]{0.49\textwidth}\centering\includegraphics[width=\linewidth,height=0.36\textheight,keepaspectratio]{folio/79_teleportation.pdf}\end{minipage}
\end{minipage}\par\bigskip\medskip
\begin{minipage}{\textwidth}
\textbf{80\quad Black-box circuit}\hfill\texttt{80\_black\_box\_circuit.py}\par
{\small Black-box quantization in four steps for the device of example 72: (a) the transmon coupled to its readout resonator; (b) the junction split into its linear inductance and its nonlinear remainder; (c) the linear part as one admittance, probed by a swept current source; (d) that admittance beside the nonlinear element, whose zero-point phase each mode sets (computed).}\par\smallskip
\begin{minipage}[t]{0.49\textwidth}\centering\includegraphics[width=\linewidth,height=0.36\textheight,keepaspectratio]{field/80_black_box_circuit.pdf}\end{minipage}\hfill
\begin{minipage}[t]{0.49\textwidth}\centering\includegraphics[width=\linewidth,height=0.36\textheight,keepaspectratio]{folio/80_black_box_circuit.pdf}\end{minipage}
\end{minipage}\par\bigskip\medskip
\begin{minipage}{\textwidth}
\textbf{81\quad Recipe grid}\hfill\texttt{81\_recipe\_grid.py}\par
{\small A generic transmon design recipe in eight numbered steps: each box holds instruction lines, a mini-plot from toy models and the tools the step uses, arrows run through the grid and two optional steps sit in dashed boxes below.}\par\smallskip
\begin{minipage}[t]{0.49\textwidth}\centering\includegraphics[width=\linewidth,height=0.36\textheight,keepaspectratio]{field/81_recipe_grid.pdf}\end{minipage}\hfill
\begin{minipage}[t]{0.49\textwidth}\centering\includegraphics[width=\linewidth,height=0.36\textheight,keepaspectratio]{folio/81_recipe_grid.pdf}\end{minipage}
\end{minipage}\par\bigskip\medskip
\begin{minipage}{\textwidth}
\textbf{82\quad Step guide}\hfill\texttt{82\_step\_guide.py}\par
{\small A generic validation guide in two columns: the numbered steps run down each column with a mini-plot beside the text and the tools as labelled chips, and two optional steps sit in dashed side boxes joined to their step.}\par\smallskip
\begin{minipage}[t]{0.49\textwidth}\centering\includegraphics[width=\linewidth,height=0.36\textheight,keepaspectratio]{field/82_step_guide.pdf}\end{minipage}\hfill
\begin{minipage}[t]{0.49\textwidth}\centering\includegraphics[width=\linewidth,height=0.36\textheight,keepaspectratio]{folio/82_step_guide.pdf}\end{minipage}
\end{minipage}\par\bigskip\medskip
\begin{minipage}{\textwidth}
\textbf{83\quad Hamiltonian terms}\hfill\texttt{83\_hamiltonian\_terms.py}\par
{\small A generic transmon-resonator Hamiltonian: (a) every term on the tile of the circuit element that produces it, bracketed by origin, (b) the one-excitation block as a bracketed matrix with its bare and dressed frequencies, (c) the drive term and the inputs that stay assumed (illustrative numbers).}\par\smallskip
\begin{minipage}[t]{0.49\textwidth}\centering\includegraphics[width=\linewidth,height=0.36\textheight,keepaspectratio]{field/83_hamiltonian_terms.pdf}\end{minipage}\hfill
\begin{minipage}[t]{0.49\textwidth}\centering\includegraphics[width=\linewidth,height=0.36\textheight,keepaspectratio]{folio/83_hamiltonian_terms.pdf}\end{minipage}
\end{minipage}\par\bigskip\medskip
\begin{minipage}{\textwidth}
\textbf{84\quad Design flow}\hfill\texttt{84\_design\_flow.py}\par
{\small A generic chip design flow on a three-by-three grid: each box lists its quantities with a status tag (specified, assumed, simulated, derived, measured), arrows snake through the steps and a dashed arrow carries the measured frequency back to the junction (illustrative numbers).}\par\smallskip
\begin{minipage}[t]{0.49\textwidth}\centering\includegraphics[width=\linewidth,height=0.36\textheight,keepaspectratio]{field/84_design_flow.pdf}\end{minipage}\hfill
\begin{minipage}[t]{0.49\textwidth}\centering\includegraphics[width=\linewidth,height=0.36\textheight,keepaspectratio]{folio/84_design_flow.pdf}\end{minipage}
\end{minipage}\par\bigskip\medskip
\begin{minipage}{\textwidth}
\textbf{85\quad Screening flow}\hfill\texttt{85\_screening\_flow.py}\par
{\small A generic literature-screening flow chart: reports pass a screening box, the exclusions branch off a trunk, a dotted arrow marks the citation search that adds reports to the full-text set, and the kept reports split into labelled categories (made-up counts).}\par\smallskip
\begin{minipage}[t]{0.49\textwidth}\centering\includegraphics[width=\linewidth,height=0.36\textheight,keepaspectratio]{field/85_screening_flow.pdf}\end{minipage}\hfill
\begin{minipage}[t]{0.49\textwidth}\centering\includegraphics[width=\linewidth,height=0.36\textheight,keepaspectratio]{folio/85_screening_flow.pdf}\end{minipage}
\end{minipage}\par\bigskip\medskip
\begin{minipage}{\textwidth}
\textbf{86\quad Mask cell with magnifiers}\hfill\texttt{86\_mask\_zoom.py}\par
{\small Xmon cell drawn to scale: a claw couples the quarter-wave resonator, a mask-level SQUID bridges the bottom gap, and two circular magnifiers show the claw and the SQUID, still to scale (designed geometry).}\par\smallskip
\begin{minipage}[t]{0.49\textwidth}\centering\includegraphics[width=\linewidth,height=0.36\textheight,keepaspectratio]{field/86_mask_zoom.pdf}\end{minipage}\hfill
\begin{minipage}[t]{0.49\textwidth}\centering\includegraphics[width=\linewidth,height=0.36\textheight,keepaspectratio]{folio/86_mask_zoom.pdf}\end{minipage}
\end{minipage}\par\bigskip\medskip
\begin{minipage}{\textwidth}
\textbf{87\quad Bowtie sensor}\hfill\texttt{87\_bowtie\_sensor.py}\par
{\small Resonator-free bowtie sensor: (a) the two wedge islands to scale beside a feedline, with a charge line below and a schematic magnifier on the junction; (b) its circuit, with the gate set by a dc source through C\_g (designed geometry).}\par\smallskip
\begin{minipage}[t]{0.49\textwidth}\centering\includegraphics[width=\linewidth,height=0.36\textheight,keepaspectratio]{field/87_bowtie_sensor.pdf}\end{minipage}\hfill
\begin{minipage}[t]{0.49\textwidth}\centering\includegraphics[width=\linewidth,height=0.36\textheight,keepaspectratio]{folio/87_bowtie_sensor.pdf}\end{minipage}
\end{minipage}\par\bigskip\medskip
\begin{minipage}{\textwidth}
\textbf{88\quad Cavity plan}\hfill\texttt{88\_cavity\_plan.py}\par
{\small Two-chamber 3D cavity in plan view, to scale: a storage and a readout chamber joined by a trench that holds a transmon chip with two antenna pads, a coupling pin through the readout wall, and a uniform field washed over both chambers (designed geometry).}\par\smallskip
\begin{minipage}[t]{0.49\textwidth}\centering\includegraphics[width=\linewidth,height=0.36\textheight,keepaspectratio]{field/88_cavity_plan.pdf}\end{minipage}\hfill
\begin{minipage}[t]{0.49\textwidth}\centering\includegraphics[width=\linewidth,height=0.36\textheight,keepaspectratio]{folio/88_cavity_plan.pdf}\end{minipage}
\end{minipage}\par\bigskip\medskip
\begin{minipage}{\textwidth}
\textbf{89\quad Topologies and a coloured circuit}\hfill\texttt{89\_topology\_protocol.py}\par
{\small (a) Coupling graphs of five chip topologies: two qubits without a coupler, a chain, a star, a shared bus and a 2 x 3 lattice. (b) The four-qubit chain's sensing circuit, each wire in its qubit's colour: GHZ encoding, a shaded exposure window where a field acts on every wire at once, decoding and readout.}\par\smallskip
\begin{minipage}[t]{0.49\textwidth}\centering\includegraphics[width=\linewidth,height=0.36\textheight,keepaspectratio]{field/89_topology_protocol.pdf}\end{minipage}\hfill
\begin{minipage}[t]{0.49\textwidth}\centering\includegraphics[width=\linewidth,height=0.36\textheight,keepaspectratio]{folio/89_topology_protocol.pdf}\end{minipage}
\end{minipage}\par\bigskip\medskip
\section*{\color{scqPurple}Comparison}
\begin{minipage}{\textwidth}
\textbf{90\quad Stacked budgets}\hfill\texttt{90\_stacked\_budget.py}\par
{\small Dark-count budget of four photon counters split by source with each total printed beside its bar, and the share of four value classes in three study groups (synthetic numbers).}\par\smallskip
\begin{minipage}[t]{0.49\textwidth}\centering\includegraphics[width=\linewidth,height=0.36\textheight,keepaspectratio]{field/90_stacked_budget.pdf}\end{minipage}\hfill
\begin{minipage}[t]{0.49\textwidth}\centering\includegraphics[width=\linewidth,height=0.36\textheight,keepaspectratio]{folio/90_stacked_budget.pdf}\end{minipage}
\end{minipage}\par\bigskip\medskip
\begin{minipage}{\textwidth}
\textbf{91\quad Efficiency chain}\hfill\texttt{91\_efficiency\_chain.py}\par
{\small Detection efficiency of three photon counters as a product of four factors on a log axis with the measured efficiency marked, and parameter values per report stacked by status and grouped by study type (synthetic numbers).}\par\smallskip
\begin{minipage}[t]{0.49\textwidth}\centering\includegraphics[width=\linewidth,height=0.36\textheight,keepaspectratio]{field/91_efficiency_chain.pdf}\end{minipage}\hfill
\begin{minipage}[t]{0.49\textwidth}\centering\includegraphics[width=\linewidth,height=0.36\textheight,keepaspectratio]{folio/91_efficiency_chain.pdf}\end{minipage}
\end{minipage}\par\bigskip\medskip
\begin{minipage}{\textwidth}
\textbf{92\quad Factor budget}\hfill\texttt{92\_factor\_bars.py}\par
{\small Signal of two devices against a reference design as signed log factors with the net change, and one device's budget as bars coloured by sign (synthetic numbers).}\par\smallskip
\begin{minipage}[t]{0.49\textwidth}\centering\includegraphics[width=\linewidth,height=0.36\textheight,keepaspectratio]{field/92_factor_bars.pdf}\end{minipage}\hfill
\begin{minipage}[t]{0.49\textwidth}\centering\includegraphics[width=\linewidth,height=0.36\textheight,keepaspectratio]{folio/92_factor_bars.pdf}\end{minipage}
\end{minipage}\par\bigskip\medskip
\begin{minipage}{\textwidth}
\textbf{93\quad Grouped bars}\hfill\texttt{93\_grouped\_bars.py}\par
{\small Logarithmic sensitivities of a transmon's frequency, anharmonicity, coupling and dispersive shift to four design parameters (computed), and surface participation of two designs with the change printed over each pair (synthetic numbers).}\par\smallskip
\begin{minipage}[t]{0.49\textwidth}\centering\includegraphics[width=\linewidth,height=0.36\textheight,keepaspectratio]{field/93_grouped_bars.pdf}\end{minipage}\hfill
\begin{minipage}[t]{0.49\textwidth}\centering\includegraphics[width=\linewidth,height=0.36\textheight,keepaspectratio]{folio/93_grouped_bars.pdf}\end{minipage}
\end{minipage}\par\bigskip\medskip
\begin{minipage}{\textwidth}
\textbf{94\quad Tornado}\hfill\texttt{94\_tornado.py}\par
{\small One-at-a-time sensitivity of a projected coupling reach to six design inputs, each moved to the ends of its range, with the range and the design step that sets each input (computed from a power-law model).}\par\smallskip
\begin{minipage}[t]{0.49\textwidth}\centering\includegraphics[width=\linewidth,height=0.36\textheight,keepaspectratio]{field/94_tornado.pdf}\end{minipage}\hfill
\begin{minipage}[t]{0.49\textwidth}\centering\includegraphics[width=\linewidth,height=0.36\textheight,keepaspectratio]{folio/94_tornado.pdf}\end{minipage}
\end{minipage}\par\bigskip\medskip
\begin{minipage}{\textwidth}
\textbf{95\quad Bar rows}\hfill\texttt{95\_bar\_rows.py}\par
{\small Background probability per measurement in two studies, with a hatched upper limit, open simulated terms and the values printed beside the bars, and measured reduction factors drawn out from 1 with published values as grey ticks (synthetic numbers).}\par\smallskip
\begin{minipage}[t]{0.49\textwidth}\centering\includegraphics[width=\linewidth,height=0.36\textheight,keepaspectratio]{field/95_bar_rows.pdf}\end{minipage}\hfill
\begin{minipage}[t]{0.49\textwidth}\centering\includegraphics[width=\linewidth,height=0.36\textheight,keepaspectratio]{folio/95_bar_rows.pdf}\end{minipage}
\end{minipage}\par\bigskip\medskip
\begin{minipage}{\textwidth}
\textbf{96\quad Dot rows and bounds}\hfill\texttt{96\_dot\_rows.py}\par
{\small Predicted (open) against measured (filled) efficiency per detector, energy thresholds per sensor as points and ranges with where each energy is counted, and efficiency against dark-count rate with arrows on the two bounds (synthetic numbers).}\par\smallskip
\begin{minipage}[t]{0.49\textwidth}\centering\includegraphics[width=\linewidth,height=0.36\textheight,keepaspectratio]{field/96_dot_rows.pdf}\end{minipage}\hfill
\begin{minipage}[t]{0.49\textwidth}\centering\includegraphics[width=\linewidth,height=0.36\textheight,keepaspectratio]{folio/96_dot_rows.pdf}\end{minipage}
\end{minipage}\par\bigskip\medskip
\begin{minipage}{\textwidth}
\textbf{97\quad Study strips}\hfill\texttt{97\_study\_strips.py}\par
{\small Four quantities across six devices on strips that share the device rows, each on its own log axis; filled markers are measured, open markers simulated, and n.r. marks a quantity a source does not report (synthetic numbers).}\par\smallskip
\begin{minipage}[t]{0.49\textwidth}\centering\includegraphics[width=\linewidth,height=0.36\textheight,keepaspectratio]{field/97_study_strips.pdf}\end{minipage}\hfill
\begin{minipage}[t]{0.49\textwidth}\centering\includegraphics[width=\linewidth,height=0.36\textheight,keepaspectratio]{folio/97_study_strips.pdf}\end{minipage}
\end{minipage}\par\bigskip\medskip
\begin{minipage}{\textwidth}
\textbf{98\quad Ratio strip and stems}\hfill\texttt{98\_ratio\_strip.py}\par
{\small Predicted over measured values by parameter type with each point's status and published scatters as grey bars, and forecast (open) against measured lifetime per device joined by a stem, grouped by chip (synthetic numbers).}\par\smallskip
\begin{minipage}[t]{0.49\textwidth}\centering\includegraphics[width=\linewidth,height=0.36\textheight,keepaspectratio]{field/98_ratio_strip.pdf}\end{minipage}\hfill
\begin{minipage}[t]{0.49\textwidth}\centering\includegraphics[width=\linewidth,height=0.36\textheight,keepaspectratio]{folio/98_ratio_strip.pdf}\end{minipage}
\end{minipage}\par\bigskip\medskip
\begin{minipage}{\textwidth}
\textbf{99\quad Timeline}\hfill\texttt{99\_timeline.py}\par
{\small Reports by year in five research lines, one marker per report shaped by its kind; the sparse years before 2018 share a compressed stretch of the axis, and reports of one line and year sit side by side (synthetic record).}\par\smallskip
\begin{minipage}[t]{0.49\textwidth}\centering\includegraphics[width=\linewidth,height=0.36\textheight,keepaspectratio]{field/99_timeline.pdf}\end{minipage}\hfill
\begin{minipage}[t]{0.49\textwidth}\centering\includegraphics[width=\linewidth,height=0.36\textheight,keepaspectratio]{folio/99_timeline.pdf}\end{minipage}
\end{minipage}\par\bigskip\medskip
\end{document}
